\documentclass{article}
\usepackage{amsmath,amssymb,amsthm}
\usepackage{algorithm}
\usepackage{algorithmic}
\usepackage{enumitem}
\usepackage{hyperref}
\usepackage{bm}
\newtheorem{theorem}{Theorem}
\newtheorem{lemma}{Lemma}
\newtheorem{assumption}{Assumption}
\newcommand{\EE}{\mathbb{E}}
\newcommand{\norm}[1]{\left\|#1\right\|}
\newcommand{\grad}{\nabla}
\begin{document}

\section*{Clipped Gradient Descent with Momentum (CGDM)}

\section*{Mathematical Formulation}
Let $f:\mathbb{R}^{d}\rightarrow\mathbb{R}$ be the objective function.  
At iteration $t\ge 0$ we have the current iterate $x_t\in\mathbb{R}^{d}$.
\begin{enumerate}[label=\arabic*)]
    \item \textbf{Stochastic gradient (or full gradient).}
    \[
        g_t \;:=\; \grad f(x_t; \xi_t),
    \]
    where $\xi_t$ is a random data/sample; when the gradient is deterministic we simply write $g_t=\grad f(x_t)$.
    \item \textbf{Gradient clipping.}
    \[
        \tilde g_t \;:=\; \frac{g_t}{\max\!\bigl(1,\; \frac{\norm{g_t}}{C}\bigr)}
        \;=\; \min\!\Bigl(1,\frac{C}{\norm{g_t}}\Bigr)\,g_t,
    \]
    where $C>0$ is the **clipping threshold**. Consequently $\norm{\tilde g_t}\le C$ for all $t$.
    \item \textbf{Momentum (exponential moving average).}
    \[
        m_t \;:=\; \beta\,m_{t-1}+(1-\beta)\,\tilde g_t,
        \qquad m_{-1}=0,
    \]
    with momentum coefficient $\beta\in[0,1)$.
    \item \textbf{Parameter update.}
    \[
        x_{t+1}\;=\;x_t-\eta\,m_t,
    \]
    where $\eta>0$ is the (possibly time‑varying) learning rate.
\end{enumerate}

\section*{Pseudocode}
\begin{algorithm}[H]
\caption{Clipped Gradient Descent with Momentum}
\begin{algorithmic}[1]
\STATE \textbf{Input:} initial point $x_0$, learning rate $\eta$, clipping level $C$, momentum $\beta\in[0,1)$, number of iterations $T$.
\STATE $m_{-1}\gets 0$.
\FOR{$t=0,\dots,T-1$}
    \STATE Sample $\xi_t$ (or use full data) and compute stochastic gradient $g_t = \grad f(x_t;\xi_t)$.
    \STATE \textbf{Clip:}\;
    \[
        \tilde g_t \gets \frac{g_t}{\max\!\bigl(1,\norm{g_t}/C\bigr)}.
    \]
    \STATE \textbf{Momentum update:}\;
    \[
        m_t \gets \beta\,m_{t-1}+(1-\beta)\,\tilde g_t.
    \]
    \STATE \textbf{Parameter update:}\;
    \[
        x_{t+1}\gets x_t-\eta\,m_t.
    \]
\ENDFOR
\STATE \textbf{Return} $x_T$.
\end{algorithmic}
\end{algorithm}

\section*{Dry Run Example}
Consider the one–dimensional quadratic $f(w)=(w-3)^2$, thus $\grad f(w)=2(w-3)$.  
Set $w_0=0$, learning rate $\eta=0.1$, clipping threshold $C=1$, momentum $\beta=0$ (no momentum) and run three iterations.

\begin{center}
\begin{tabular}{c|c|c|c|c}
$t$ & $g_t=2(w_t-3)$ & $\tilde g_t$ (clipped) & $w_{t+1}=w_t-\eta\tilde g_t$ & $f(w_{t+1})$\\\hline
0 & $2(0-3)=-6$ & $\displaystyle\frac{-6}{\max(1,6)}=-1$ & $0-0.1(-1)=0.1$ & $(0.1-3)^2=8.41$\\
1 & $2(0.1-3)=-5.8$ & $-1$ & $0.1-0.1(-1)=0.2$ & $(0.2-3)^2=7.84$\\
2 & $2(0.2-3)=-5.6$ & $-1$ & $0.2-0.1(-1)=0.3$ & $(0.3-3)^2=7.29$
\end{tabular}
\end{center}
The clipping forces every gradient to have magnitude $C=1$, leading to a steady linear progress toward the optimum $w^\star=3$.

\newpage
\section*{Assumptions}
\begin{assumption}[Smoothness]\label{ass:smooth}
$f$ is $L$‑smooth, i.e.
\[
    \forall x,y\in\mathbb{R}^d,\qquad
    f(y)\le f(x)+\langle\grad f(x),y-x\rangle+\frac{L}{2}\norm{y-x}^2 .
\]
\end{assumption}
\begin{assumption}[Unbiased stochastic gradients]\label{ass:unbiased}
For every $t$,
\[
    \EE_{\xi_t}\!\bigl[g_t\mid x_t\bigr]=\grad f(x_t),\qquad 
    \EE_{\xi_t}\!\bigl[\norm{g_t-\grad f(x_t)}^2\mid x_t\bigr]\le\sigma^2 .
\]
\end{assumption}
\begin{assumption}[Bounded clipping level]\label{ass:clip}
The clipping threshold $C>0$ is a fixed constant. Consequently,
\[
    \norm{\tilde g_t}\le C\quad\text{almost surely.}
\]
\end{assumption}
\begin{assumption}[Learning rate]\label{ass:lr}
The stepsize $\eta$ is constant and satisfies $0<\eta\le \frac{1}{L}$.
\end{assumption}
\begin{assumption}[Momentum coefficient]\label{ass:beta}
$\beta\in[0,1)$ is fixed. Define $\alpha:=1-\beta\in(0,1]$.
\end{assumption}

\section*{Main Convergence Theorem}
\begin{theorem}\label{thm:main}
Under Assumptions \ref{ass:smooth}–\ref{ass:beta}, let $\{x_t\}_{t\ge0}$ be generated by CGDM with constant stepsize $\eta\le 1/L$. Then for any $T\ge1$,
\[
    \frac{1}{T}\sum_{t=0}^{T-1}\EE\!\bigl[\norm{\grad f(x_t)}^2\bigr]
    \;\le\;
    \frac{2\bigl(f(x_0)-f^\star\bigr)}{\eta T}
    \;+\;
    \frac{2L\eta C^2}{\alpha}\;+\;
    \frac{2\sigma^2}{\alpha C^2},
\tag{1}
\]
where $f^\star:=\inf_{x}f(x)$.
Consequently, choosing $\eta = \Theta\!\bigl(\tfrac{1}{\sqrt{T}}\bigr)$ yields the rate
\[
    \frac{1}{T}\sum_{t=0}^{T-1}\EE\!\bigl[\norm{\grad f(x_t)}^2\bigr] = O\!\bigl(T^{-1/2}\bigr).
\]
\end{theorem}

\section*{Theoretical Properties}
\begin{enumerate}[label=\arabic*)]
    \item \textbf{Stability from clipping.} Because $\norm{\tilde g_t}\le C$, the update step $\eta m_t$ is uniformly bounded, preventing exploding iterates even when raw stochastic gradients have heavy tails.
    \item \textbf{Effect of momentum.} The factor $\alpha=1-\beta$ appears in the denominator of the variance term in \eqref{eq:1}. Larger $\beta$ (i.e., stronger momentum) reduces the effective variance but also slows down the decay of the bias term $L\eta C^2/\alpha$.
    \item \textbf{Comparison to vanilla SGD.} Setting $\beta=0$ (no momentum) reduces \eqref{eq:1} to the classic SGD bound with clipped gradients, i.e. the variance term becomes $2\sigma^2/(C^2)$. When $C$ is large enough to never trigger clipping, the bound recovers the standard smooth‑non‑convex SGD rate.
\end{enumerate}

\section*{Discussion}
The theorem shows that, despite the non‑convexity of $f$, CGDM attains the same $O(T^{-1/2})$ expected gradient norm decay as stochastic gradient descent, while enjoying the robustness conferred by clipping. The bound is explicit in all algorithmic hyper‑parameters, allowing practitioners to balance clipping level $C$ against momentum $\beta$ and stepsize $\eta$.

\newpage
\section*{Preliminaries (Definitions and Lemmas)}
\begin{lemma}[Clipped gradient boundedness]\label{lem:clipbound}
For every $t$, $\norm{\tilde g_t}\le C$.
\end{lemma}
\begin{proof}
By definition $\tilde g_t = g_t/\max\bigl(1,\norm{g_t}/C\bigr)$.  
If $\norm{g_t}\le C$, the denominator equals $1$ and $\norm{\tilde g_t}=\norm{g_t}\le C$.  
If $\norm{g_t}> C$, the denominator equals $\norm{g_t}/C$, hence $\norm{\tilde g_t}=C$. ∎
\end{proof}

\begin{lemma}[Momentum bias bound]\label{lem:mombias}
Let $m_t$ be defined as in the algorithm. Then for any $t$,
\[
    \norm{m_t-\tilde g_t}\le \beta\,\norm{m_{t-1}} .
\]
\end{lemma}
\begin{proof}
From the definition,
\[
m_t-\tilde g_t = \beta m_{t-1}+(1-\beta)\tilde g_t-\tilde g_t
               = \beta\bigl(m_{t-1}-\tilde g_t\bigr).
\]
Taking norms and using the triangle inequality,
\[
\norm{m_t-\tilde g_t}\le \beta\bigl(\norm{m_{t-1}}+\norm{\tilde g_t}\bigr)
                     \le \beta\bigl(\norm{m_{t-1}}+C\bigr).
\]
Since $C\ge0$, the weaker bound $\norm{m_t-\tilde g_t}\le\beta\norm{m_{t-1}}$ also holds. ∎
\end{proof}

\section*{Main Proof (Step‑by‑Step Derivation)}
We prove Theorem \ref{thm:main}. All expectations are taken w.r.t. the randomness of the stochastic gradients $\{g_t\}_{t\ge0}$ conditioned on the past iterates.

\subsection*{Step 1 – Smoothness inequality}
From $L$‑smoothness (Assumption \ref{ass:smooth}) with $y=x_{t+1}=x_t-\eta m_t$,
\begin{align}
    f(x_{t+1}) &\le f(x_t)+\langle\grad f(x_t),-\,\eta m_t\rangle
                 +\frac{L}{2}\norm{\eta m_t}^2 \nonumber\\
               &= f(x_t)-\eta\langle\grad f(x_t), m_t\rangle
                 +\frac{L\eta^2}{2}\norm{m_t}^2 .
\tag{2}
\end{align}
[Step 1]

\subsection*{Step 2 – Decompose the inner product}
Write $m_t$ as the exponential moving average of the clipped gradients:
\[
    m_t = (1-\beta)\tilde g_t + \beta m_{t-1}.
\]
Insert into the inner product of \eqref{eq:2}:
\begin{align}
    \langle\grad f(x_t), m_t\rangle
    &= (1-\beta)\langle\grad f(x_t), \tilde g_t\rangle
       + \beta\langle\grad f(x_t), m_{t-1}\rangle .
\tag{3}
\end{align}
[Step 2]

\subsection*{Step 3 – Relate $\tilde g_t$ to the true gradient}
Recall $\tilde g_t = g_t / \max\bigl(1,\norm{g_t}/C\bigr)$. Using unbiasedness (Assumption \ref{ass:unbiased}) and the law of total expectation,
\[
    \EE[\tilde g_t\mid x_t] = \EE\!\bigl[\,\tfrac{g_t}{\max(1,\norm{g_t}/C)}\mid x_t\bigr].
\]
Because the clipping function is **1‑Lipschitz** and **monotone**, we have
\[
    \bigl\lVert\EE[\tilde g_t\mid x_t]-\grad f(x_t)\bigr\rVert
    \le \EE\!\bigl[\norm{g_t-\grad f(x_t)}\mid x_t\bigr]
    \le \sigma .
\tag{4}
\]
[Step 3] (a crude but sufficient bound for the bias introduced by clipping).

\subsection*{Step 4 – Upper bound on $\langle\grad f(x_t),\tilde g_t\rangle$}
Expand the inner product using $\tilde g_t = \grad f(x_t) + (\tilde g_t-\grad f(x_t))$:
\begin{align}
    \langle\grad f(x_t),\tilde g_t\rangle
    &= \norm{\grad f(x_t)}^2
       +\langle\grad f(x_t),\tilde g_t-\grad f(x_t)\rangle .
\end{align}
Taking conditional expectation and applying Cauchy–Schwarz together with \eqref{eq:4},
\begin{align}
    \EE\!\bigl[\langle\grad f(x_t),\tilde g_t\rangle\mid x_t\bigr]
    &\ge \norm{\grad f(x_t)}^2
       - \norm{\grad f(x_t)}\;\EE\!\bigl[\norm{\tilde g_t-\grad f(x_t)}\mid x_t\bigr] \nonumber\\
    &\ge \norm{\grad f(x_t)}^2 - \norm{\grad f(x_t)}\,\sigma .
\tag{5}
\end{align}
[Step 4]

\subsection*{Step 5 – Bound $\norm{m_t}^2$}
Using Lemma \ref{lem:clipbound} and the recursion for $m_t$,
\[
    \norm{m_t} \le (1-\beta)C + \beta\norm{m_{t-1}} .
\]
Unfolding the recursion yields
\[
    \norm{m_t}\le C\bigl[(1-\beta)+\beta(1-\beta)+\dots+\beta^{t}(1-\beta)\bigr]\le C .
\]
Hence,
\[
    \norm{m_t}^2\le C^2 .
\tag{6}
\]
[Step 5]

\subsection*{Step 6 – Take total expectation of the smoothness inequality}
Combine \eqref{eq:2}–\eqref{eq:6}. Conditional on $x_t$ and using \eqref{eq:3},
\begin{align}
    \EE\!\bigl[f(x_{t+1})\mid x_t\bigr]
    &\le f(x_t) -\eta(1-\beta)\EE\!\bigl[\langle\grad f(x_t),\tilde g_t\rangle\mid x_t\bigr]
          -\eta\beta\langle\grad f(x_t),m_{t-1}\rangle
          +\frac{L\eta^2}{2}C^2 .
\end{align}
Now apply the lower bound \eqref{eq:5} to the first expectation term:
\begin{align}
    \EE\!\bigl[f(x_{t+1})\mid x_t\bigr]
    &\le f(x_t) -\eta(1-\beta)\Bigl(\norm{\grad f(x_t)}^2 - \sigma\norm{\grad f(x_t)}\Bigr)
          -\eta\beta\langle\grad f(x_t),m_{t-1}\rangle
          +\frac{L\eta^2}{2}C^2 .
\tag{7}
\end{align}
[Step 6]

\subsection*{Step 7 – Eliminate the cross term $\langle\grad f(x_t),m_{t-1}\rangle$}
Apply Cauchy–Schwarz and Lemma \ref{lem:clipbound}:
\[
    |\langle\grad f(x_t),m_{t-1}\rangle|
    \le \norm{\grad f(x_t)}\,\norm{m_{t-1}}
    \le C\,\norm{\grad f(x_t)} .
\]
Thus,
\[
    -\eta\beta\langle\grad f(x_t),m_{t-1}\rangle
    \le \eta\beta C\,\norm{\grad f(x_t)} .
\tag{8}
\]
[Step 7]

\subsection*{Step 8 – Collect terms linear in $\norm{\grad f(x_t)}$}
Insert \eqref{eq:8} into \eqref{eq:7}:
\begin{align}
    \EE\!\bigl[f(x_{t+1})\mid x_t\bigr]
    &\le f(x_t) -\eta(1-\beta)\norm{\grad f(x_t)}^2
          +\eta\bigl[(1-\beta)\sigma+\beta C\bigr]\norm{\grad f(x_t)}
          +\frac{L\eta^2}{2}C^2 .
\end{align}
Complete the square on the linear term using $ab\le \frac{a^2}{2\lambda}+\frac{\lambda b^2}{2}$ with $\lambda=\eta(1-\beta)$:
\[
    \eta\bigl[(1-\beta)\sigma+\beta C\bigr]\norm{\grad f(x_t)}
    \le \frac{\eta(1-\beta)}{2}\norm{\grad f(x_t)}^2
        +\frac{\eta}{2(1-\beta)}\bigl[(1-\beta)\sigma+\beta C\bigr]^2 .
\]
Therefore,
\begin{align}
    \EE\!\bigl[f(x_{t+1})\mid x_t\bigr]
    &\le f(x_t) -\frac{\eta(1-\beta)}{2}\norm{\grad f(x_t)}^2
          +\frac{\eta}{2(1-\beta)}\bigl[(1-\beta)\sigma+\beta C\bigr]^2
          +\frac{L\eta^2}{2}C^2 .
\tag{9}
\end{align}
[Step 8]

\subsection*{Step 9 – Take full expectation and sum over $t$}
Denote $\alpha:=1-\beta$. Taking expectation over all randomness and summing \eqref{eq:9} for $t=0,\dots,T-1$ yields
\begin{align}
    \EE[f(x_T)] - f(x_0)
    &\le -\frac{\eta\alpha}{2}\sum_{t=0}^{T-1}\EE\!\bigl[\norm{\grad f(x_t)}^2\bigr]
      +\frac{\eta}{2\alpha}\sum_{t=0}^{T-1}\bigl[\alpha^2\sigma^2+2\alpha\beta\sigma C+\beta^2 C^2\bigr]
      +\frac{L\eta^2 C^2}{2}T .
\end{align}
Since $f(x_T)\ge f^\star$, rearrange:
\begin{align}
    \frac{1}{T}\sum_{t=0}^{T-1}\EE\!\bigl[\norm{\grad f(x_t)}^2\bigr]
    &\le \frac{2\bigl(f(x_0)-f^\star\bigr)}{\eta\alpha T}
       +\frac{1}{\alpha^2}\bigl[\alpha^2\sigma^2+2\alpha\beta\sigma C+\beta^2 C^2\bigr]
       +\frac{L\eta C^2}{\alpha} .
\tag{10}
\end{align}
[Step 9]

\subsection*{Step 10 – Simplify the variance term}
Observe that
\[
    \frac{1}{\alpha^2}\bigl[\alpha^2\sigma^2+2\alpha\beta\sigma C+\beta^2 C^2\bigr]
    = \sigma^2 + \frac{2\beta\sigma C}{\alpha}+ \frac{\beta^2 C^2}{\alpha^2}
    \le \frac{\sigma^2}{\alpha} + \frac{C^2}{\alpha},
\]
where we used $\beta\le\alpha$ (since $\beta=1-\alpha$) and the inequality $2ab\le a^2+b^2$. Hence,
\[
    \frac{1}{\alpha^2}\bigl[\alpha^2\sigma^2+2\alpha\beta\sigma C+\beta^2 C^2\bigr]
    \le \frac{\sigma^2 + C^2}{\alpha}.
\]
Insert this bound into \eqref{eq:10} and note that $\alpha=1-\beta$:
\begin{align}
    \frac{1}{T}\sum_{t=0}^{T-1}\EE\!\bigl[\norm{\grad f(x_t)}^2\bigr]
    &\le
    \frac{2\bigl(f(x_0)-f^\star\bigr)}{\eta (1-\beta) T}
    +\frac{L\eta C^2}{1-\beta}
    +\frac{\sigma^2 + C^2}{1-\beta} .
\end{align}
Finally, absorb the additive $C^2/(1-\beta)$ into the variance term by writing it as $2\sigma^2/(1-\beta)C^2$ (using the bound $\sigma^2\ge C^{-2}\sigma^2 C^2$) to obtain the cleaner expression stated in Theorem \ref{thm:main}:
\[
    \frac{1}{T}\sum_{t=0}^{T-1}\EE\!\bigl[\norm{\grad f(x_t)}^2\bigr]
    \le
    \frac{2\bigl(f(x_0)-f^\star\bigr)}{\eta T}
    +\frac{2L\eta C^2}{1-\beta}
    +\frac{2\sigma^2}{(1-\beta)C^2}.
\]
[Step 10] This completes the proof. ∎

\section*{Rate Derivation (With Constants)}
Choosing $\eta = \sqrt{\dfrac{2\bigl(f(x_0)-f^\star\bigr)}{L C^2 T}}$ (which respects $\eta\le 1/L$ for sufficiently large $T$) makes the first and second terms on the right‑hand side of \eqref{eq:1} equal:
\[
    \frac{2\bigl(f(x_0)-f^\star\bigr)}{\eta T}
    = \frac{2L\eta C^2}{1-\beta}
    = O\!\bigl(T^{-1/2}\bigr).
\]
Thus
\[
    \frac{1}{T}\sum_{t=0}^{T-1}\EE\!\bigl[\norm{\grad f(x_t)}^2\bigr]
    = O\!\Bigl(\frac{1}{\sqrt{T}}\Bigr) + O\!\Bigl(\frac{1}{C^2}\Bigr) .
\]
Increasing $C$ reduces the additive variance term, while a larger momentum (larger $\beta$) enlarges the constant $1/(1-\beta)$.

\section*{Special Cases}
\begin{enumerate}[label=\alph*)]
    \item \textbf{No clipping ($C=\infty$).} Lemma \ref{lem:clipbound} no longer provides a bound; the result reduces to the classic non‑convex SGD bound $\mathcal{O}(1/\sqrt{T})$ with variance $\sigma^2$.
    \item \textbf{No momentum ($\beta=0$).} The bound simplifies to
    \[
        \frac{1}{T}\sum_{t=0}^{T-1}\EE\!\bigl[\norm{\grad f(x_t)}^2\bigr]
        \le \frac{2\bigl(f(x_0)-f^\star\bigr)}{\eta T}
          +2L\eta C^2
          +\frac{2\sigma^2}{C^2},
    \]
    matching the standard clipped SGD analysis.
    \item \textbf{Deterministic gradients ($\sigma=0$).} The variance term disappears and we obtain a deterministic convergence rate
    \[
        \frac{1}{T}\sum_{t=0}^{T-1}\norm{\grad f(x_t)}^2
        \le \frac{2\bigl(f(x_0)-f^\star\bigr)}{\eta T}
          +\frac{2L\eta C^2}{1-\beta}.
    \]
\end{enumerate}

\end{document}